\section{State of Research}

The piezoresistive effect has been extensively investigated in both experimental and engineering contexts, 
particularly in semiconductor-based sensors and microelectromechanical systems (MEMS) \cite{smith1954,doll2013}. 
Early studies established the strong dependence of electrical conductivity on mechanical strain in materials such as silicon, 
highlighting the importance of crystal structure and carrier transport in determining the sensitivity of the effect.

From a materials physics perspective, piezoresistivity originates from strain-induced modifications of the electronic structure 
and charge transport properties. Mechanical deformation can alter band structure, effective masses, and scattering mechanisms, 
which in turn influence carrier mobility and electrical conductivity \cite{hamaguchi}. 
These effects are well understood at the microscopic and transport levels.

In practical modeling and engineering applications, however, this coupling is typically introduced in a simplified manner, 
where the conductivity is treated as an empirical function of strain. While such approaches are effective for device design, 
they do not provide a transparent link between microscopic transport mechanisms and macroscopic behavior.

At the continuum level, partial differential equation models are widely used to describe both elasticity and electrical conduction. 
Nevertheless, the incorporation of strain-dependent transport properties into these models is usually carried out phenomenologically, 
without a systematic connection to the underlying physics.

This reveals a gap between physically grounded transport descriptions and continuum-level modeling. 
Bridging this gap in a computationally consistent way remains an open and relevant problem, particularly for materials and devices 
where electromechanical coupling plays a central role.

Based on these considerations, the proposed study is guided by the following research questions:

\begin{itemize}
    \item How can strain-dependent conductivity be incorporated into a consistent continuum model of electromechanical coupling?
    \item How does mechanical deformation influence the spatial distribution of electrical potential in conductive materials?
    \item What effects arise from the nonlinear coupling between mechanics and electrical transport?
    \item To what extent can flexible (e.g.\ data-driven or physics-informed) approaches improve the modeling of $\sigma(\varepsilon)$?
\end{itemize}

To address these questions, the study is structured around the following components:

\begin{itemize}
    \item Physical background and modeling assumptions
    \item Formulation of the coupled electromechanical model
    \item Analysis of the nonlinear system
    \item Computational methods and numerical experiments
    \item Interpretation of results in the context of material behavior
\end{itemize}